\chapter{eSim Installation and Setup}

\section{System Requirements}

Development and testing for the OpenROAD integration were performed on \textbf{Ubuntu 22.04 LTS}. The following system configuration was used during implementation and testing:

\begin{itemize}
    \item \textbf{Operating System:} Ubuntu 22.04 LTS
    \item \textbf{RAM:} Minimum 4 GB (8 GB recommended)
    \item \textbf{Storage:} At least 20 GB free disk space
    \item \textbf{Internet Connection:} Required for repository cloning and dependency installation
    \item \textbf{Python Version:} Python 3.8 or later with PyQt5 support
\end{itemize}

The installation environment also included OpenROAD-flow-scripts (ORFS), KLayout, Yosys, and additional Python dependencies required for backend automation.

\section{Standard eSim Installation}

The eSim source code was cloned from the official FOSSEE GitHub repository and installed using the provided installation script.

\begin{lstlisting}[language=bash, caption={Standard eSim installation commands.}]
$ git clone https://github.com/FOSSEE/eSim.git
$ cd eSim
$ chmod +x install-eSim.sh
$ ./install-eSim.sh --install
\end{lstlisting}

The installation script automatically installs the required packages including Python, PyQt5, Ngspice, KiCad, and related system dependencies.

\section{Required Dependencies}

Additional dependencies were installed to support OpenROAD integration and RTL-to-GDSII backend flow execution.

\begin{lstlisting}[language=bash, caption={Additional dependencies for OpenROAD integration.}]
$ sudo apt update

$ sudo apt install -y build-essential git cmake
$ sudo apt install -y python3-dev python3-pip
$ sudo apt install -y klayout

$ pip3 install --user pyqt5
\end{lstlisting}

These dependencies were required for GUI execution, Python backend scripting, layout visualization, and ORFS workflow execution.

\section{Workspace Setup and Project Creation}

After installation, eSim was launched from the terminal or application menu. New projects were created using the \textit{New Project} option available in the eSim graphical interface.

Each project generated its own directory containing schematic files, generated netlists, simulation outputs, and backend design files required during OpenROAD execution.

\section{GUI Overview and OpenROAD Button}

The eSim graphical interface provides project management, schematic editing, simulation control, and tool integration support. During this internship, an additional \textbf{OpenROAD} execution button was integrated into the main toolbar of the application.

The button was connected with the backend automation flow implemented in \texttt{OpenROAD.py}. When triggered, it automatically processed the selected project and initiated the OpenROAD RTL-to-GDSII flow.

Figure~\ref{fig:openroad_button} shows the OpenROAD integration button added to the eSim GUI toolbar.

\begin{figure}[H]
\centering
\includegraphics[width=0.50\textwidth]{OpenROAD-Button.png}
\caption{OpenROAD integration button added to the eSim GUI toolbar.}
\label{fig:openroad_button}
\end{figure}

\section{Installation Verification}

After installation and integration, the setup was verified by launching eSim and confirming that:

\begin{itemize}
    \item the graphical interface opened correctly,
    \item the OpenROAD button was visible in the toolbar,
    \item existing projects could be loaded successfully,
    \item generated \texttt{.cir.out} files were valid,
    \item OpenROAD Flow Scripts executed without errors.
\end{itemize}

Figure~\ref{fig:terminal_install} shows the terminal output generated during project execution and ORFS backend flow processing.

\begin{figure}[H]
\centering
\includegraphics[width=0.92\textwidth]{FA-Terminal-process.png}
\caption{Terminal output showing the eSim project build and ORFS execution process.}
\label{fig:terminal_install}
\end{figure}